% !TeX root = ../../tesis.tex
%%%%%%%%%%%%%%%%%%%%%%%%%%%%%%%%%%%%%%%%%%%%%%%%%%%%%%%%%%%%%%%%%%%
% SECCIÓN 6.3 — ESCENARIOS E INDICADORES DE LA RED CON ANILLO
% Issue #25.
%
% Esquema común de §6.3.2–§6.3.5:
%   1. Entrada: qué mide el indicador (definición en Cap. 3) y su
%      valor en el Caso 1 (Cap. 5).
%   2. \subsubsection{Caso 2: Anillo MB}    — valores y mapa
%   3. \subsubsection{Caso 3: Anillo Metro} — valores y mapa
%      (si ambos casos dan lo mismo: una sola \subsubsection
%      «Casos 2 y 3» con nota de la diferencia)
%   4. Lectura comparativa + figura comparativa, si existe.
% Figuras: se deciden caso por caso (sin tope fijo).
%%%%%%%%%%%%%%%%%%%%%%%%%%%%%%%%%%%%%%%%%%%%%%%%%%%%%%%%%%%%%%%%%%%
\section{Escenarios e Indicadores de la Red con Anillo}
\label{sec:anillo-indicadores}

% §6.3.1 — Los escenarios del API-I y sus razones
% !TeX root = ../../tesis.tex
%%%%%%%%%%%%%%%%%%%%%%%%%%%%%%%%%%%%%%%%%%%%%%%%%%%%%%%%%%%%%%%%%%%
% §6.3.1 — LOS ESCENARIOS DEL API-I
%
% Estructura:
%   · Entrada: por qué dos casos (mismo trazo, dos modos de operación)
%   · Caso 2: Anillo MB — razones (las da el autor), vía confinada,
%     mapa Propuesta_Anillar_Extendido_MB
%   · Caso 3: Anillo Metro — razones (las da el autor), vía exclusiva,
%     mapa Propuesta_Anillar_Extendido_Metro (trasladado desde §6.1.3)
%   · Cierre: cuadro comparativo de los casos en términos urbanos,
%     sin adelantar resultados
% Las notas de cuadrante de los mapas extendidos traen resultados
% (B(v), T): no usarlos aquí; van en §6.3.2–§6.3.5.
%%%%%%%%%%%%%%%%%%%%%%%%%%%%%%%%%%%%%%%%%%%%%%%%%%%%%%%%%%%%%%%%%%%

% Figura trasladada desde §6.1.3 (2026-10-18); se integra al
% redactar el Caso 3.
\begin{figure}[H]
    \centering
    \includegraphics[width=0.92\textwidth]{%
        Figures/Mapas/Qgis_Resources/RedConAnillo/Escenario_Metro/Base/Propuesta_Anillar_Extendido_Metro}
    \caption[Cuadrantes del corredor anillar]{%
    Corredor anillar propuesto por cuadrante (norponiente, nororiente,
    surponiente y suroriente) y sistemas de transporte que cruza.
    Corrida: 2026-10-04.}
    \label{fig:propuesta-anillo-cuadrantes}
    \fuentefigura{Elaboración propia con \textbf{QGIS};
    datos: \texttt{Apímetro} y \texttt{VFTModel}~\citep{cabrera2025apimetro, cabrera2026VFTModel}.}
\end{figure}


% §6.3.2 — Nivel de Cobertura Espacial (C)
% !TeX root = ../../tesis.tex
%%%%%%%%%%%%%%%%%%%%%%%%%%%%%%%%%%%%%%%%%%%%%%%%%%%%%%%%%%%%%%%%%%%
% §6.3.1 — NIVEL DE COBERTURA ESPACIAL (C) EN LOS CASOS CON ANILLO
% (archivo 6-3-2 por la numeración anterior). Redactada 2026-10-22.
%
% Datos (corrida 2026-10-04; RESPUESTAS_PREGUNTAS_ISSUES.md:160,532 y
%   garibelt/comparativo/cobertura_comparativa.csv):
%   C ZMVM 76 = 13.29 % → 13.43 % (0.166 → 0.168); CDMX 53.52 → 53.53 %;
%   municipios 3.92 → 4.09 %; 42 → 41 municipios en cero (Texcoco);
%   7 demarcaciones cambian; Casos 2 y 3 idénticos (mismas estaciones).
%   Figuras usan dominio 141 (0.055 → 0.056), decisión C6-1.
% PENDIENTE #34: recalcular el ámbito CDMX (16 alcaldías) al final.
% Figuras (decisión del autor: las dos en el cuerpo y en el Anexo G
%   a página completa):
%   Clasificacion_Garibelt_Mapa_Cobertura_MB (= _Metro)
%   Comparativo_Garibelt_Cobertura
%%%%%%%%%%%%%%%%%%%%%%%%%%%%%%%%%%%%%%%%%%%%%%%%%%%%%%%%%%%%%%%%%%%
\subsection{Nivel de Cobertura Espacial (\texorpdfstring{$C$}{C})}
\label{subsec:anillo-cobertura}

La cobertura indica qué parte del territorio queda a una distancia caminable
de alguna estación: es el área dentro de \termino{buffers} de 800~m alrededor
de cada parada, como proporción de la superficie de cada demarcación
(§\ref{subsec:indicador-cobertura}). En el Caso~1 la red cubre el 13.29\,\% de
la \zmvm{}, un promedio que reúne situaciones muy distintas. Más de la mitad
de la \cdmx{} (53.52\,\%) está cerca de una estación; en los 60 municipios
conurbados la proporción baja a 3.92\,\% y 42 de ellos no tienen cobertura
(§\ref{subsec:cobertura-espacial}). Normalizado, el indicador vale 0.166 y
queda en el intervalo 0--0.25 del Espectro MJG
(Cuadro~\ref{tab:espectro-mjg-diagnostico}).

\subsubsection{Casos~2 y~3}

En la cobertura los dos casos con anillo dan el mismo resultado. El indicador
depende solo de dónde están las estaciones, no de la velocidad ni del derecho
de vía, y el Anillo MB y el Anillo Metro tienen las mismas 94. Con ellas, la
superficie cubierta crece unos 11~km$^2$ y la cobertura de la \zmvm{} pasa de
13.29 a 13.43\,\%. En la \cdmx{} queda prácticamente igual (53.53\,\%) y en
los municipios conurbados sube a 4.09\,\%. El valor normalizado pasa de 0.166
a 0.168, dentro del mismo intervalo. La
Figura~\ref{fig:anillo-cobertura-mb} muestra la cobertura por demarcación con
el anillo; el mapa del Caso~2 vale también para el Caso~3, porque los dos son
idénticos en este indicador.

\begin{figure}[H]
    \centering
    \includegraphics[width=\textwidth]{%
        Figures/Mapas/Qgis_Resources/RedConAnillo/Escenario_MB/Clasificacion_Garibelt/Clasificacion_Garibelt_Mapa_Cobertura_MB}
    \caption[Cobertura por demarcación en la red con anillo]{%
    Cobertura espacial por demarcación en el Caso~2 (Anillo MB), igual a la
    del Caso~3. El mapa usa las 141 demarcaciones del conjunto de datos; las
    cifras del texto corresponden a las 76 de la \zmvm{}. Versión a página
    completa en la Figura~\ref{fig:anx-g-cobertura-mb}. Corrida: 2026-10-04.}
    \label{fig:anillo-cobertura-mb}
    \fuentefigura{Elaboración propia con \textbf{QGIS};
    datos: \texttt{Apímetro} y \texttt{VFTModel}~\citep{cabrera2025apimetro, cabrera2026VFTModel}.}
\end{figure}

Como se observa en la Figura~\ref{fig:anillo-cobertura-comparativo}, siete
demarcaciones registran algún cambio. Solo en dos se nota: Nezahualcóyotl, que
pasa de 39.8 a 48.6\,\%, y Tlalnepantla de Baz, que pasa de 53.9 a 58.9\,\%;
las dos están en el Estado de México, sobre el norte y el oriente del anillo.
Ecatepec de Morelos, Texcoco, Naucalpan de Juárez, Iztapalapa y Xochimilco
ganan menos de medio punto, y ninguna demarcación cambia de rango en el mapa.
Con el anillo, Texcoco deja la cobertura cero y los municipios sin cobertura
bajan de 42 a 41.

\begin{figure}[H]
    \centering
    \includegraphics[width=\textwidth]{%
        Figures/Mapas/Qgis_Resources/Comparativo_CG/Comparativo_Garibelt_Cobertura}
    \caption[Cobertura por demarcación: red actual y red con anillo]{%
    Cobertura espacial por demarcación en el Caso~1 (A) y en los Casos~2 y~3
    (B), que coinciden porque comparten estaciones. Las demarcaciones
    achuradas son las de mayor ganancia. El mapa usa las 141 demarcaciones del
    conjunto de datos; las cifras del texto corresponden a las 76 de la
    \zmvm{}. Versión a página completa en la
    Figura~\ref{fig:anx-g-cobertura-comparativo}. Corrida: 2026-10-04.}
    \label{fig:anillo-cobertura-comparativo}
    \fuentefigura{Elaboración propia con \textbf{QGIS};
    datos: \texttt{Apímetro} y \texttt{VFTModel}~\citep{cabrera2025apimetro, cabrera2026VFTModel}.}
\end{figure}

\subsubsection{Lectura Comparativa}

El anillo casi no amplía la cobertura, y la razón está en su trazo. El
Periférico recorre la parte de la ciudad que ya concentra más estaciones, así
que en la \cdmx{} los \termino{buffers} de las estaciones nuevas caen casi
siempre sobre área que ya estaba cubierta. Donde el anillo sale a los
municipios, en el norte y el oriente, sí agrega superficie, porque ahí la red
estructurada es escasa. La cobertura mide cercanía, no conexión. Un anillo que
acerca pocas estaciones a territorio nuevo puede, aun así, cambiar los
recorridos dentro de la red, y eso es lo que revisan la Fuerza Capilar, el
tiempo de viaje y la intermediación en los apartados siguientes. Los 41
municipios que siguen sin cobertura están fuera del alcance del \textsc{api-i}
y son el territorio que tendría que atender un anillo exterior
(§\ref{subsec:corredor-anillar-limitaciones}).


% §6.3.3 — Fuerza Capilar (FC)
% !TeX root = ../../tesis.tex
%%%%%%%%%%%%%%%%%%%%%%%%%%%%%%%%%%%%%%%%%%%%%%%%%%%%%%%%%%%%%%%%%%%
% §6.3.3 — FUERZA CAPILAR (FC) EN LOS CASOS CON ANILLO
%
% Datos (corrida 2026-10-04): mediana 7 (Q1 6, Q3 10) → 0.250 Débil
%   en los tres casos; entre MB y Metro cambia el FC de 13 nodos,
%   todos sobre el Periférico; 2 bajan de Débil a Crítico (Alameda
%   Oriente, Muyuguarda/Cafetales). Fuente: nota del mapa FC Metro.
% Estructura: entrada · Casos 2 y 3 (posible subsubsección única
%   con nota de diferencias) · lectura comparativa.
% Figuras candidatas:
%   .../Escenario_MB/Clasificacion_Garibelt/Clasificacion_Garibelt_Fuerza_Capilar_MB
%   .../Escenario_Metro/Clasificacion_Garibelt/Clasificacion_Garibelt_Fuerza_Capilar_Metro
%%%%%%%%%%%%%%%%%%%%%%%%%%%%%%%%%%%%%%%%%%%%%%%%%%%%%%%%%%%%%%%%%%%


% §6.3.4 — Tiempo de Viaje Promedio (T)
% !TeX root = ../../tesis.tex
%%%%%%%%%%%%%%%%%%%%%%%%%%%%%%%%%%%%%%%%%%%%%%%%%%%%%%%%%%%%%%%%%%%
% §6.3.4 — TIEMPO DE VIAJE PROMEDIO (T) EN LOS CASOS CON ANILLO
%
% Datos (corrida 2026-10-04): T = 113.91 (Caso 1, 0.6957 Aceptable)
%   → 108.99 (Caso 2, −4.92 min, −4.3 %, 0.7475 Aceptable, a 0.0025
%   de Idóneo) → 101.92 (Caso 3, −11.99 min, −10.5 %, 0.8219 Idóneo).
%   Tabla 3.5 del Cap. 3 anticipaba −24 % (suposición).
%   Referencia: «tiempo de referencia del modelo» (85 min, D7).
% Estructura: entrada · Caso 2 · Caso 3 · lectura comparativa.
% Figuras: no hay mapa de T; cuadro LaTeX.
%%%%%%%%%%%%%%%%%%%%%%%%%%%%%%%%%%%%%%%%%%%%%%%%%%%%%%%%%%%%%%%%%%%


% §6.3.5 — Centralidad de Intermediación B(v)
% !TeX root = ../../tesis.tex
%%%%%%%%%%%%%%%%%%%%%%%%%%%%%%%%%%%%%%%%%%%%%%%%%%%%%%%%%%%%%%%%%%%
% §6.3.5 — CENTRALIDAD DE INTERMEDIACIÓN B(v) EN LOS CASOS CON ANILLO
%
% Datos (corrida 2026-10-04; bv_comparativo.csv):
%   Tacubaya 0.2232 → 0.2163 (Caso 2, −3.1 %) → 0.1569 (Caso 3, −29.7 %)
%   Caso 2: ninguna estación del anillo en el top 50; MB conserva los
%     10 nodos de la red actual (Periférico Oriente/Tláhuac 0.061, #69)
%   Caso 3: 9/10 y 40/50 del top son del anillo; top 5 = Tacubaya,
%     Periférico Oriente/Tláhuac 0.1550, Rómulo O'Farril 0.1310,
%     Barranca del Muerto/Periférico 0.1284, Constitución de
%     1917/Periférico 0.1274; 208 nodos pierden centralidad.
%   Gini 0.8007 → 0.8012 → 0.8134 (0.199 · 0.199 · 0.187, Crítico).
% Criterio de las figuras (decisión C6-5):
%   comparativos = orden por nodo; Dashboard_Betweenness_* = 10
%   nombres con mayor suma de B(v) (como tab:bv-top10-agregado).
% Estructura: entrada · Caso 2 · Caso 3 · lectura comparativa.
% Figuras candidatas:
%   .../Clasificacion_Garibelt_Betweenness_Centrality_MB / _Metro
%   Tableau_Resources/RedConAnillo/Escenario_MB|Metro/Dashboard_Betweenness_*
%   Qgis_Resources/Comparativo_CG/Comparativo_Garibelt_Betweenness
%   Tableau_Resources/Comparativo_CG/Dashboard_Comparativo_Intermediacion_Bv
%%%%%%%%%%%%%%%%%%%%%%%%%%%%%%%%%%%%%%%%%%%%%%%%%%%%%%%%%%%%%%%%%%%

